\pdfminorversion=7
\pdfsuppresswarningpagegroup=1
\documentclass[
  candidate, % document type
  14pt, % font size
  subf % use and configure subfig package for nested figure numbering
]{disser}

% Кодировка и язык
\usepackage[T2A]{fontenc} % поддержка кириллицы
\usepackage[utf8]{inputenc} % кодировка исходного текста
\usepackage[english,russian]{babel} % переключение языков
\usepackage{indentfirst} % ПАКЕТ, чтобы был отступ текста после заголовка
\usepackage{ragged2e}    % ДОБАВЬТЕ ЭТОТ ПАКЕТ для умного выравнивания

% Геометрия страницы и графика
\usepackage[left=30mm, right=15mm, top=20mm, bottom=20mm]{geometry} % поля страницы
\usepackage{graphicx} % подключение графики
\usepackage{pdfpages} % вставка pdf-страниц
\usepackage{float} % управление положением плавающих объектов

% Таблицы
\usepackage{array} % расширенные возможности для работы с таблицами
\usepackage{tabularx} % автоматический подбор ширины столбцов
\usepackage{dcolumn} % выравнивание чисел по разделителю
\newcolumntype{Y}{>{\centering\arraybackslash}X}

% Математика
\usepackage{bm} % полужирное начертание для математических символов
\usepackage{amsmath} % дополнительные математические возможности
\usepackage{amssymb} % дополнительные математические символы

% Библиография и ссылки
\usepackage{cite} % поддержка цитирования
\usepackage{hyperref} % создание гиперссылок
\hypersetup{hidelinks}

% Прочее
\usepackage{color} % работа с цветом
\usepackage{epstopdf} % конвертация eps в pdf
\usepackage{multirow} % объединение ячеек таблиц по вертикали
\usepackage{afterpage} % вставка материала после текущей страницы
\usepackage[font={normal},compatibility=false]{caption} % настройка подписей к рисункам и таблицам
\usepackage[onehalfspacing]{setspace} % полуторный интервал
\usepackage{fancyhdr} % установка колонтитулов
\usepackage{listings} % поддержка вставки исходного кода
\usepackage{enumitem}

\setlist{nosep} % Делает все списки максимально компактными

% Установка шрифта Times New Roman
\renewcommand{\rmdefault}{ftm}

\makeatletter
\renewcommand{\@dotsep}{1}
\providecommand{\@chapapp}{\chaptername}
\providecommand{\chaptername}{}
\@ifundefined{prethechapter}{}{\renewcommand{\prethechapter}{}}
\@ifundefined{postthechapter}{}{\renewcommand{\postthechapter}{\space}}
\@ifundefined{prethesection}{}{\renewcommand{\prethesection}{}}
\@ifundefined{postthesection}{}{\renewcommand{\postthesection}{\space}}
\@ifundefined{prethesubsection}{}{\renewcommand{\prethesubsection}{}}
\@ifundefined{postsubsection}{}{\renewcommand{\postsubsection}{\space}}
\@ifundefined{prethesubsubsection}{}{\renewcommand{\prethesubsubsection}{}}
\@ifundefined{postsubsubsection}{}{\renewcommand{\postsubsubsection}{\space}}

% 2. Префиксы и постфиксы для ОГЛАВЛЕНИЯ (ToC)
\@ifundefined{tocprethechapter}{}{\renewcommand{\tocprethechapter}{}}          % Убирает слово "Глава"
\@ifundefined{tocpostthechapter}{}{\renewcommand{\tocpostthechapter}{\space}}  % Убирает точку после номера главы
\@ifundefined{tocpostthesection}{}{\renewcommand{\tocpostthesection}{\space}}  % Убирает точку после номера раздела
\@ifundefined{tocpostthesubsection}{}{\renewcommand{\tocpostthesubsection}{\space}} % Убирает точку после номера подраздела
\@ifundefined{tocpostthesubsubsection}{}{\renewcommand{\tocpostthesubsubsection}{\space}}

% Уменьшаем ширину блока под номер в оглавлении
% Формат: \@dottedtocline{уровень}{отступ слева}{ширина блока под номер}
\renewcommand*{\l@chapter}{\@dottedtocline{1}{0pt}{0pt}}
\renewcommand*{\l@section}{\@dottedtocline{1}{1.1em}{2em}} % было 1.1em 1.5em
\renewcommand*{\l@subsection}{\@dottedtocline{2}{2.1em}{2.3em}}
% \renewcommand*{\l@subsubsection}{\@dottedtocline{3}{3.1em}{3.1em}} либо 3.5em

\newcommand{\centeredheading}[1]{%
  \par\medskip
  \begin{center}
    \normalfont\bfseries\normalsize #1
  \end{center}
  \par\medskip
}
\makeatother

% Заголовки глав и разделов: обычный размер, но жирный шрифт, выравнивание по левому краю
\renewcommand\thechapteralign{\RaggedRight}
\renewcommand\thechapterfont{\normalfont\bfseries}
\renewcommand\chapteralign{\RaggedRight}
\renewcommand\chapterfont{\normalfont\bfseries}
\renewcommand\schapteralign{\centering}
\renewcommand\schapterfont{\normalfont\bfseries}

\renewcommand\sectionfont{\normalfont\bfseries}
\renewcommand\subsectionfont{\normalfont\bfseries}
\renewcommand\subsubsectionfont{\normalfont\bfseries}
\renewcommand\paragraphfont{\normalfont\bfseries}
\renewcommand\subparagraphfont{\normalfont\bfseries}
\renewcommand\sectionalign{\RaggedRight}
\renewcommand\subsectionalign{\RaggedRight}
\renewcommand\subsubsectionalign{\RaggedRight}
\renewcommand\subparagraphalign{\RaggedRight}

\setlength{\parindent}{1.25cm} % глобальный отступ для основного текста

\renewcommand\chapterindent{1.25cm}
\renewcommand\schapterindent{0pt}
\renewcommand\sectionindent{1.25cm}
\renewcommand\subsectionindent{1.25cm}
\renewcommand\subsubsectionindent{1.25cm}
\renewcommand\paragraphindent{1.25cm}
\renewcommand\subparagraphindent{1.25cm}

% Нормальный шрифт для пунктов оглавления
\renewcommand\tocchapterfont{\normalfont}
\renewcommand\tocchapterfillfont{\normalfont}
\renewcommand\tocchapternumfont{\normalfont}
\renewcommand\tocsectionfont{\normalfont}
\renewcommand\tocsectionfillfont{\normalfont}
\renewcommand\tocsectionnumfont{\normalfont}
\renewcommand\tocsubsectionfont{\normalfont}
\renewcommand\tocsubsectionfillfont{\normalfont}
\renewcommand\tocsubsectionnumfont{\normalfont}
\renewcommand\tocsubsubsectionfont{\normalfont}
\renewcommand\tocsubsubsectionfillfont{\normalfont}
\renewcommand\tocsubsubsectionnumfont{\normalfont}

% Создание нового типа столбца для выравнивания содержимого по центру
\newcommand{\PreserveBackslash}[1]{\let\temp=\\#1\let\\=\temp}
\newcolumntype{C}[1]{>{\PreserveBackslash\centering}p{#1}}

% Настройка стиля страницы
\pagestyle{fancy}      % Использование стиля "fancy" для оформления страниц
\fancyhf{}              % Очистка текущих значений колонтитулов
\fancyfoot[C]{\thepage} % Установка номера страницы в нижнем колонтитуле по центру
\renewcommand{\headrulewidth}{0pt} % Удаление разделительной линии в верхнем колонтитуле

% Настройка подписей к изображениям и таблицам
\captionsetup{format=hang,labelsep=period}

% Использование полужирного начертания для векторов
\let\vec=\mathbf

% Установка глубины оглавления
\setcounter{tocdepth}{2}

% Указание папки для поиска изображений
\graphicspath{{images/}}

% Установка стилей страницы и главы
\pagestyle{footcenter}
\chapterpagestyle{footcenter}

% 2. Гарантированное отключение вывода ссылок \cite по всему тексту
\AtBeginDocument{
  \renewcommand{\cite}[1]{}
}

% Команда для инлайн-кода (как обратные кавычки в Markdown)
\newcommand{\code}[1]{\mbox{\texttt{#1}}}

\emergencystretch=3em % Разрешает дополнительное растяжение пробелов для предотвращения Overfull \hbox